\subsection{CP Violation from String Theory --- arXiv:1808.07060}

\textbf{Authors:} Hans Peter Nilles, Michael Ratz, Andreas Trautner, Patrick K. S. Vaudrevange

\paragraph{主な主張}
ヘテロティック $\mathbb{Z}_3$ orbifoldの $\Delta(54)$ フレーバー対称性で、軽い状態だけに定義できるCPが重い巻き付き状態との結合によって破れる具体的な機構を示した\cite{Nilles:2018wex}。

\paragraph{新規性}
CPをフレーバー群の外部自己同型として検討し、重い弦状態が幾何学的CP破れを引き起こすことを弦理論のスペクトルから実現した。

\paragraph{位置付け}
低エネルギーのCP対称性を論じる際に重い弦状態の表現も重要になることを示す研究である。

\paragraph{コメント（読書メモ）}
読書メモ：$\mathbb{Z}_3$ orbifoldではCPの候補を $S_4$ の中に探せる、と理解した。ここで $S_4$ は $\Delta(54)$ の「一つの外部自己同型」ではなく外部自己同型群 $\operatorname{Out}(\Delta(54))\simeq S_4$ であり、その元がCP候補になると表現を整える。ただし実際に存在する全ての表現に対してCPとして働くかは別の条件で、重い状態を含めた判定が必要になる。
